\chapter{Системийн туршилт, үр дүн, хэлэлцүүлэг}
\label{Chapter4}

\section{Туршилтын зорилго}
Хөгжүүлсэн системийн туршилтын зорилго нь функциональ шаардлагууд хэрэгжсэн эсэх, үндсэн хэрэглээний урсгал тасралтгүй ажиллаж байгаа эсэх, хиймэл оюун болон ахицын логик бодит өгөгдөл дээр зөв ажиллаж байгаа эсэхийг баталгаажуулахад оршино. Туршилт нь зөвхөн код ажиллаж байна уу гэдгийг шалгах бус, хэрэглэгчийн өнцгөөс системийн уялдаа холбоо, мэдээллийн зөв дамжуулалт, нөөц механизм, аюулгүй байдлын суурь шийдэл хангалттай эсэхийг дүгнэх зорилготой.

\section{Туршилтын орчин}
Туршилтыг дараах орчинд хийхээр төлөвлөсөн.
\begin{itemize}
\item Урд талын хэсэг: React + Vite хөгжүүлэлтийн орчин
\item Арын хэсэг: Django хөгжүүлэлтийн сервер
\item Өгөгдлийн сан: Docker доторх PostgreSQL эсвэл локал SQLite
\item API туршилт: хөтөч, урд-арын хэсгийн уялдаа, шаардлагатай тохиолдолд REST клиент
\item Хиймэл оюуны орчин: Gemini API түлхүүртэй болон API түлхүүргүй хоёр горим
\end{itemize}

Ийм хоёр горим ашиглах нь хиймэл оюуны үйлчилгээ хэвийн үед болон нөөц нөхцөлд систем ямар байдалтай ажиллахыг харьцуулахад чухал ач холбогдолтой.

\section{Интерфейсийн туршилтын дүрслэлүүд}
Энэ хэсэгт системийн үндсэн урсгалуудыг дэлгэцийн зураг хэлбэрээр баримтжуулан оруулав. Зураг бүр нь хэрэглэгчийн тодорхой үйлдэл, модуль хоорондын уялдаа холбоо, системийн бодит хэрэгжилтийг харуулах зорилготой.

\subsection{Нэвтрэлт ба бүртгэлийн интерфейс}
\insertimage{zurag/login.png}{Системийн нэвтрэх интерфейс}

\subsection{Түвшин тогтоох шалгалтын урсгал}
\insertimage{zurag/placement-gen.png}{Түвшин тогтоох шалгалт ажиллаж буй байдал}
\insertimage{zurag/placement-early.png}{Түвшин тогтоох шалгалтын асуултыг дахин үүсгэж буй процесс}
\insertimage{zurag/placement-result.png}{Түвшин тогтоох шалгалтын үр дүн болон түвшин тогтоосон байдал}

\subsection{Түвшин ахих шалгалтын урсгал}
\insertimage{zurag/levelup.png}{Түвшин ахих шалгалтын интерфейсийн жишээ}
\insertimage{zurag/levelup-gen.png}{Түвшин ахих шалгалт үүсгэж буй байдал}
\insertimage{zurag/levelup-result.png}{Түвшин ахих шалгалтын үр дүн}

\subsection{Курс ба суралцах замнал}
\insertimage{zurag/all-course.png}{Курсийн жагсаалтын дэлгэц}
\insertimage{zurag/course-track.png}{Суралцах замналын хуудас}

\subsection{Хичээлийн дэлгэрэнгүй хуудасны дүрслэл}
\insertimage{zurag/lessons.png}{Хичээлийн хуудасны ерөнхий бүтэц}
\insertimage{zurag/lesson-top.png}{Хичээлийн дэлгэрэнгүй хуудасны дээд хэсэг}
\insertimage{zurag/lesson-mid.png}{Хичээлийн дэлгэрэнгүй хуудасны дунд хэсгийн мэдээлэл}
\insertimage{zurag/lesson-bottom.png}{Хураангуй, сорил болон хичээлийн доод хэсэг}

\subsection{Ахицын хураангуйн дүрслэл}
\insertimage{zurag/progress-top.png}{Progress хэсгийн дээд талын мэдээлэл}
\insertimage{zurag/progress-mid.png}{Progress хэсгийн дунд талын мэдээлэл}
\insertimage{zurag/progress-bottom.png}{Progress хэсгийн доод талын аналитик мэдээлэл}

\subsection{Нэмэлт боломжууд}
\insertimage{zurag/chatbot.png}{Хиймэл оюуны чатбот ашиглаж буй интерфейс}
\insertimage{zurag/profile-dropdown.png}{Нүүр хуудасны profile цэсний харагдац}
\insertimage{zurag/profile-edit.png}{Profile хуудасны засварлах интерфейс}
\insertimage{zurag/liked-lesson.png}{Таалагдсан хичээлүүдийн хуудасны харагдац}

\subsection{Тоглоомын хэсгийн дүрслэл}
\insertimage{zurag/games.png}{Тоглоомын хэсгийн ерөнхий харагдац}
\insertimage{zurag/css-game.png}{CSS төрлийн интерактив тоглоомын жишээ}
\insertimage{zurag/html-game.png}{HTML төрлийн интерактив тоглоомын жишээ}

\section{Туршилтын аргачлал}
Системийн туршилтыг дараах төрлүүдээр ангилан авч үзэв.
\begin{enumerate}
\item Функциональ туршилт
\item Интеграцийн туршилт
\item UI урсгалын туршилт
\item Хиймэл оюуны нөөц урсгалын туршилт
\item Аюулгүй байдал ба хязгаарлалтын туршилт
\end{enumerate}

\subsection{Функциональ туршилт}
Модуль бүр өөрийн үндсэн шаардлагыг хангаж байгаа эсэхийг шалгана. Жишээлбэл бүртгэлийн endpoint шинэ хэрэглэгч үүсгэж байна уу, курсийн жагсаалт зөв шүүлттэй ирж байна уу, \texttt{complete\_lesson} дуудагдахад ахиц болон хэрэглэгчийн оноо зөв шинэчлэгдэж байна уу гэдгийг шалгана.

\subsection{Интеграцийн туршилт}
Урд ба арын хэсгийн хооронд өгөгдөл зөв дамжиж буй эсэхийг шалгана. Жишээ нь хичээлийн дэлгэрэнгүй хуудас ачаалагдахад хичээлийн мэдээлэл, дадлагын даалгавар, ахиц, хамтын орчны хэлэлцүүлэг, хураангуй зэргийг дараалалтай авч чадаж байна уу гэдгийг шалгана.

\subsection{Хиймэл оюуны нөөц урсгалын туршилт}
Gemini API түлхүүр тохируулагдсан болон тохируулаагүй хоёр нөхцөлд хураангуй, сорил, дадлагын даалгавар, хяналтын урсгалыг туршина. Энэ нь гадаад үйлчилгээ тасарсан үед системийн хэрэглээ бүрэн зогсохгүй байх шаардлагыг баталгаажуулна.

\section{Туршилтын сценарийн жишээнүүд}
\subsection{Сценари 1: Шинэ хэрэглэгчийн анхны урсгал}
\begin{enumerate}
\item Хэрэглэгч бүртгэл үүсгэнэ.
\item Нэвтэрч хяналтын самбарт орно.
\item Түвшин тогтоох шалгалт өгөх шаардлагатай эсэхийг систем тодорхойлно.
\item Түвшин тогтоох шалгалт өгсний дараа ур чадварын түвшин хадгалагдана.
\item Курс болон замналын хэсэгт хэрэглэгчийн түвшинд тохирсон агуулга харагдана.
\end{enumerate}

Хүлээгдэх үр дүн: хэрэглэгчийн \texttt{has\_placement\_test=true} болж, \texttt{skill\_level} шинэчлэгдсэн байх ёстой.

\subsection{Сценари 2: Хичээл суралцах бүрэн урсгал}
\begin{enumerate}
\item Хэрэглэгч курсийн дэлгэрэнгүй хуудсаас хичээл сонгоно.
\item Хичээлийн дэлгэрэнгүй хуудас хичээлийн агуулга, видео, дадлагын даалгаврыг ачаална.
\item Хэрэглэгч хураангуй авч, сорил өгнө.
\item Дадлагын даалгаврын хэсэгт хариулт бичиж, шаардлагатай бол урьдчилан харах эсвэл ажиллуулагч ашиглана.
\item Хичээл дуусгалт хийж, ахиц шинэчлэгдэнэ.
\end{enumerate}

Хүлээгдэх үр дүн: \texttt{LessonProgress.is\_completed=true}, \texttt{completed\_at} утга хадгалагдах, \texttt{User.total\_score} болон \texttt{completed\_lessons} өсөх ёстой.

\subsection{Сценари 3: Хамтын орчны хяналт}
\begin{enumerate}
\item Хэрэглэгч зохистой агуулгатай пост бичнэ.
\item Систем хяналт хийж \texttt{allow} эсвэл \texttt{review} шийдвэр гаргана.
\item Дүрэм зөрчсөн түлхүүр үг бүхий пост оруулахыг оролдоно.
\item Систем \texttt{reject} шийдвэр гаргаж, анхааруулгын тоог нэмэгдүүлнэ.
\end{enumerate}

Хүлээгдэх үр дүн: зохистой агуулга нийтлэгдэх, дүрэм зөрчсөн агуулга хадгалагдахгүй байх.

\subsection{Сценари 4: Хиймэл оюуны нөөц урсгал}
\begin{enumerate}
\item Gemini API түлхүүргүйгээр хичээлийн хураангуй дуудах
\item Gemini API түлхүүргүйгээр дадлагын даалгавар үүсгэх шаардлагатай хичээл нээх
\item Хиймэл оюуны хяналт API түлхүүргүй нөхцөлд сэтгэгдэл оруулах
\end{enumerate}

Хүлээгдэх үр дүн: хураангуйн хэсэгт энгийн нөөц тайлбар ирэх, дадлагын даалгавар анхны бүтэцтэй үүсэх, хяналтын систем \texttt{allow} нөөц шийдэл ашиглах.

\section{Туршилтын үр дүнг бүртгэх загвар}
\begin{longtable}{|p{0.8cm}|p{4.2cm}|p{3.5cm}|p{4.5cm}|}
\hline
\textbf{\#} & \textbf{Туршилтын нэр} & \textbf{Оролт/нөхцөл} & \textbf{Хүлээгдэх үр дүн} \\
\hline
1 & Бүртгэл үүсгэх & Шинэ username, password & User амжилттай үүснэ \\
\hline
2 & Нэвтрэх & Зөв нэвтрэх мэдээлэл & JWT token авна \\
\hline
3 & Түвшин тогтоох шалгалт & Хариултын багц илгээх & Ур чадварын түвшин хадгалагдана \\
\hline
4 & Курсийн шүүлт & \texttt{my\_level=1} & Түвшинд тохирсон курсүүд ирнэ \\
\hline
5 & Хичээлийн хураангуй & Агуулгатай хичээл & Хиймэл оюун эсвэл нөөц хураангуй ирнэ \\
\hline
6 & Хичээлийн сорил & Хичээлийн агуулга & 3 асуулттай сорил ирнэ \\
\hline
7 & Хичээл дуусгах & Score илгээх & Ахиц, нийт оноо шинэчлэгдэнэ \\
\hline
8 & Хамтын орчны хяналт & Хориглосон үг бүхий текст & Татгалзсан хариу ирнэ \\
\hline
9 & Ахицын хураангуй & Нэвтэрсэн хэрэглэгч & Долоо хоногийн идэвх, курсын ахиц ирнэ \\
\hline
10 & Docker эхлэлт & Compose run & Database ба backend хэвийн асна \\
\hline
\end{longtable}

\section{Үр дүнгийн шинжилгээ}
\subsection{Шаталсан сургалтын урсгал бодитоор хэрэгжсэн}
Систем нь хэрэглэгчийн анхны түвшин тогтоох, түүнд тохирсон курс шүүх, хичээл дуусгалтаар ахиц бүртгэх, дараагийн түвшин рүү ахих боломж бүрдүүлэх гэсэн цогц урсгалыг бий болгосон. Энэ нь нэг хэрэглэгчид нэг ижил зам санал болгодог уламжлалт хандлагыг халж, суралцах замыг хувь хүний өгөгдөлтэй уялдуулсан.

\subsection{Хиймэл оюун сургалтын туршлагад бодит нэмэр өгсөн}
Хураангуй нь урт хичээлийг богино хугацаанд ойлгоход тусална. Сорил нь ойлголтоо шууд шалгах боломж олгоно. Дадлагын даалгавар нь онолын агуулгыг шууд хэрэглээнд шилжүүлэх алхам болдог. Чатбот болон хяналт нь зөвхөн агуулга бус харилцааны чанарыг ч дэмжиж байна.

\subsection{Ахицын хураангуй нь суралцах зан төлөвийг ил тод болгосон}
Дуусгасан хичээлүүд, нийт оноо, курсын ахиц, долоо хоногийн идэвх зэрэг үзүүлэлт нь хэрэглэгч өөрийн ахицыг хэмжих боломж олгодог. Энэ төрлийн ил тод байдал нь хэрэглэгчийн мотивацид эерэг нөлөөтэй.

\subsection{Fallback шийдэл нь найдвартай байдлыг нэмэгдүүлсэн}
Хиймэл оюуны үйлчилгээнээс бүрэн хамаарахгүй байхаар хураангуй болон дадлагын даалгаварт нөөц логик оруулсан нь системийг илүү уян хатан болгосон. Энэ нь инженерчлэлийн хувьд чухал шийдэл юм.

\section{Давуу тал, хязгаарлалт, хэлэлцүүлэг}
\subsection{Давуу тал}
\begin{itemize}
\item Цогц веб системийн түвшинд иж бүрэн хэрэгжилттэй
\item Шаталсан сургалтын урсгалтай
\item Хураангуй, сорил, дадлагын даалгавар, хяналт зэрэг олон талт хиймэл оюуны хэрэглээтэй
\item Ахиц хянах хэсэг нь хэрэглэгчийн ахицыг бодитой харуулдаг
\item Хамтын орчин болон тоглоомын хэсэг нь хэрэглэгчийн оролцоог нэмэгдүүлдэг
\item Docker болон орчны тохиргоотой тул өргөтгөх, нэвтрүүлэх бэлтгэлтэй
\end{itemize}

\subsection{Хязгаарлалт}
\begin{itemize}
\item Хиймэл оюунаар үүсгэсэн агуулгын чанар бүх тохиолдолд ижил түвшинд биш
\item Зөвлөмжийн логик нь дүрэмд суурилсан тул илүү нарийн хувьчлах боломж хязгаарлагдмал
\item Түвшин тогтоох болон түвшин ахих шалгалтын хэмжүүрийг илүү судалгааны үндэстэй нарийвчлах шаардлагатай
\item PDF текст гарган авах, хуулбар шалгах, илүү баялаг шинжилгээ зэрэг боломж одоогоор ороогүй
\item Автомат тестийн багц болон гүйцэтгэлийн шалгуурыг цаашид өргөжүүлэх шаардлагатай
\end{itemize}

\subsection{Хэлэлцүүлэг}
Энэхүү төслийн үр дүнгээс харахад дасан зохицох сургалт болон хиймэл оюуны тусламжийг нэгтгэсэн сургалтын платформыг дипломын ажлын түвшинд бодитоор хэрэгжүүлэх боломжтой байна. Гэхдээ чанарын дараагийн шатанд хүрэхийн тулд зөвлөмжийн логик, шалгалтын хэмжүүр, автомат туршилт, нэвтрүүлэлтийн дамжуулах шугам зэрэг хэсгийг улам боловсронгуй болгох шаардлагатай.

\section{Цаашдын хөгжүүлэлтийн санал}
\begin{enumerate}
\item Багш болон mentor-д зориулсан агуулга үүсгэх, хянах самбар хөгжүүлэх
\item Зөвлөмжийн логикт хэрэглэгчийн бодит зан төлөв, сонирхол, алдааны хэв маягийг оруулах
\item Сорилын чанарыг сайжруулахын тулд хүндрэлийн шошго, сэдвийн шошго, хариултын шинжилгээ нэмэх
\item Хамтын орчны нэр хүндийн систем болон хяналтын шалгах дараалал бий болгох
\item Гар утсанд нийцсэн болон PWA хувилбар хөгжүүлэх
\item Хиймэл оюуны загварын хариуг бүртгэж, чанарын шинжилгээ хийх
\item Нэгж тест, уялдааны тест, CI/CD дамжуулах шугам нэмж автоматжуулах
\end{enumerate}

\section{Бүлгийн дүгнэлт}
Энэ бүлэгт системийн туршилтын зорилго, аргачлал, сценарийн жишээ, интерфейсийн дүрслэл, үр дүнгийн шинжилгээ, давуу тал, хязгаарлалт, цаашдын хөгжүүлэлтийн чиглэлийг авч үзлээ. Ийнхүү төслийн туршилт ба үр дүнгийн хэсгийг нэг бүлэгт төвлөрүүлснээр тайлангийн бүтэц илүү цэгцтэй болж, системийн бодит хэрэгжилтээс гарсан дүгнэлтүүд ойлгомжтой болсон.
